\section{Quality Supervision and Sampling}
\label{app:q-architecture}

The auxiliary quality branch predicts class-specific localisation quality at each detection location, with continuous targets defined below. Each scale used a separately parameterised Q head (Table~\ref{tab:section32-1}), without normalisation layers. Its final convolution was zero-initialised before M3 training, and detached inputs and targets restricted Q-loss updates to the quality heads.

\begin{table}[htbp]
\centering
\small
\setlength{\tabcolsep}{4pt}
\renewcommand{\arraystretch}{1.15}
\begin{tabularx}{\linewidth}{@{}>{\raggedright\arraybackslash}p{0.32\linewidth}>{\raggedright\arraybackslash}X@{}}
\toprule
\textbf{Component} & \textbf{Transformation} \\
\midrule
Classification-feature projection & 384 input channels; \(1\times1\) convolution to 64 channels; SiLU \\
Regression-feature projection & 96 input channels; \(1\times1\) convolution to 64 channels; SiLU \\
Raw-distance projection & Four input channels; \(1\times1\) convolution to 16 channels; SiLU \\
Feature fusion & Concatenate 144 channels; \(3\times3\) convolution to 128 channels; SiLU \\
Quality output & \(1\times1\) convolution to six quality logits \\
\bottomrule
\end{tabularx}
\caption{Shared auxiliary quality-head architecture at each detection scale.}
\label{tab:section32-1}
\end{table}

Q eligibility followed Table~\ref{tab:q-source-class-eligibility}; P cells without eligible locations supplied no supervision.

For image n, location j and class c, the detached target was the maximum quality over retained same-class GT boxes g. With u denoting clipped IoU, predicted and reference centres \(\hat{\boldsymbol c}\) and c, and reference dimensions w and h,

\[
\begin{gathered}
q^*_{njc}=
\begin{cases}
\displaystyle
\operatorname{clip}_{[0,1]}\left\{
\max_{g\in\mathcal G_{nc}}
\left[
(2^{u_{njg}}-1)
\exp\left(
-\frac{\|\hat{\boldsymbol c}_{nj}-\boldsymbol c_g\|_2^2}
{\max(w_g^2+h_g^2,\varepsilon)}
\right)
\right]\right\},&\mathcal G_{nc}\ne\varnothing,\\[4pt]
0,&\mathcal G_{nc}=\varnothing,
\end{cases}\\[3pt]
\varepsilon=10^{-12}.
\end{gathered}
\]

IoU used non-negative areas and a union floor of \(\varepsilon\). This diagonal-relative quality target differs from the AXIS centre term.

Ordinary locations with maximum same-class IoU below 0.5 were sampled from intervals [0, 0.05), [0.05, 0.20), [0.20, 0.40), [0.40, 0.70) and [0.70, 1]. Each image–class–interval contributed at most 32 locations. The lowest interval prioritised original scores; other intervals used CPU permutations seeded by 42 + e for each image at epoch e, independently of the run seed.

For each GT box, eligible good locations were traversed through three lists: up to 11 by quality, up to 11 by original logits, then cyclically shifted native indices. Rankings were stable; the shift was 42 + e modulo the good-location count. First occurrences filled a quota of 32 unique locations, after which native positives were added by union.

Each non-empty ordinary interval contributed unit total weight; each GT selection added unit total weight. Contributions accumulated for shared locations. For accumulated weights \(\omega\), quality logits z and actual batch size B, weighted BCE with continuous targets gave

\[
\mathcal L_Q=
\frac{1}{B}\sum_{n=1}^{B}
\frac{
\sum_{j,c}\omega_{njc}\,
\ell_{\mathrm{BCE}}(z^Q_{njc},q^*_{njc})
}{\max\left(\sum_{j,c}\omega_{njc},1\right)},
\]

\FloatBarrier
